% answers_short.tex —— 2014 年·江苏卷 · 简短答案速查
% 由 tools/make_short_answers.py 自 answers.tex 的 \gksolhead/\gkanswer 提取，
% 并与 2014/2014_数学真题_江苏_答案要点.md 逐题比对后生成，勿手改（改源 MD 后重跑）。

\begin{gknotice}
\textbf{使用说明}：本节只给最终答案便于快速核对，完整推导见解析版。答案为编写者独立推导，\textbf{非官方标准答案}。
\end{gknotice}

\gkpart{填空题}
\begin{itemize}[leftmargin=2.6em]
  \item[\textbf{1.}] $\{-1,3\}$。两集合的公共元素为 $-1$ 与 $3$。
  \item[\textbf{2.}] $21$。$(5+2\mathrm i)^2=21+20\mathrm i$。
  \item[\textbf{3.}] $5$。$n$ 依次取 $1,2,3,4$ 时 $2^n\leqslant 20$，取 $5$ 时 $2^5=32>20$ 首次成立。
  \item[\textbf{4.}] $\dfrac13$。乘积为 $6$ 的取法只有 $\{1,6\}$、$\{2,3\}$ 共 $2$ 种，任取 $2$ 数共 $\mathrm C_4^2=6$ 种。
  \item[\textbf{5.}] $\dfrac\pi6$。由 $\sin\left(\dfrac{2\pi}3+\varphi\right)=\dfrac12$ 及 $0\leqslant\varphi<\pi$ 得 $\dfrac{2\pi}3+\varphi=\dfrac{5\pi}6$。
  \item[\textbf{6.}] $24$。前两组频率为 $(0.015+0.025)\times10=0.4$，故 $60\times0.4=24$ 株。
  \item[\textbf{7.}] $4$。设公比为 $q>0$，由 $a_2=1$ 得 $q^6=q^4+2q^2$，即 $q^2=2$，故 $a_6=q^4=4$。
  \item[\textbf{8.}] $\dfrac32$。由 $\dfrac{S_1}{S_2}=\dfrac94$ 得 $\dfrac{r_1}{r_2}=\dfrac32$，侧面积相等给出 $\dfrac{h_1}{h_2}=\dfrac23$，故 $\dfrac{V_1}{V_2}=\dfrac{S_1h_1}{S_2h_2}=\dfrac32$。
  \item[\textbf{9.}] $\dfrac{2\sqrt{55}}5$。圆心 $(2,-1)$ 到直线距离为 $\dfrac{3}{\sqrt5}$，弦长为 $2\sqrt{4-\dfrac95}$。
  \item[\textbf{10.}] $\left(-\dfrac{\sqrt2}2,0\right)$。只需端点函数值小于零：$f(m)=2m^2-1<0$ 且 $f(m+1)=2m^2+3m<0$。
  \item[\textbf{11.}] $-3$。由过点得 $8a+b=-10$，由切线斜率 $-7$ 得 $16a-b=-14$，解得 $a=-1$，$b=-2$。
  \item[\textbf{12.}] $22$。$\overrightarrow{AP}=\dfrac14\overrightarrow{AB}+\overrightarrow{AD}$，$\overrightarrow{BP}=\overrightarrow{AD}-\dfrac34\overrightarrow{AB}$，代入数量积条件即得。
  \item[\textbf{13.}] $\left(0,\dfrac12\right)$。一个周期内 $f(x)=a$ 有 $4$ 个根，区间 $[-3,4]$ 含两个整周期及 $[3,4]$ 上另外 $2$ 个根，共 $10$ 个根当且仅当 $0<a<\dfrac12$。
  \item[\textbf{14.}] $\dfrac{\sqrt6-\sqrt2}4$。由正弦定理 $a+\sqrt2b=2c$，令 $t=\dfrac ab$，则 $\cos C=\dfrac{3t^2-2\sqrt2t+2}{8t}$ 在 $t=\dfrac{\sqrt6}3$ 处取最小值。
\end{itemize}

\gkpart{解答题}
\begin{itemize}[leftmargin=2.6em]
  \item[\textbf{15.}] (1) $-\dfrac{\sqrt{10}}{10}$。由 $\sin\alpha=\dfrac{\sqrt5}5$，$\alpha\in\left(\dfrac\pi2,\pi\right)$ 得 $\cos\alpha=-\dfrac{2\sqrt5}5$，再展开 $\sin\left(\dfrac\pi4+\alpha\right)$；(2) $-\dfrac{4+3\sqrt3}{10}$。由 $\sin2\alpha=-\dfrac45$，$\cos2\alpha=\dfrac35$ 展开 $\cos\left(\dfrac{5\pi}6-2\alpha\right)$
  \item[\textbf{16.}] (1) 见解析版；(2) 见解析版
  \item[\textbf{17.}] (1) $\dfrac{x^2}2+y^2=1$。$BF_2=\sqrt{b^2+c^2}=a=\sqrt2$，代入 $C\left(\dfrac43,\dfrac13\right)$ 得 $b^2=1$；(2) $\dfrac{\sqrt5}5$。由 $A\left(\dfrac{2a^2c}{a^2+c^2},\dfrac{b(c^2-a^2)}{a^2+c^2}\right)$ 与 $C$ 关于 $x$ 轴对称，$\overrightarrow{F_1C}\cdot\overrightarrow{AB}=0$ 化为 $5a^2c^2=a^4$
  \item[\textbf{18.}] (1) $150\ \mathrm m$。以 $O$ 为原点建系得 $B(80,120)$，$BC=\sqrt{90^2+120^2}=150$；(2) $OM=10\ \mathrm m$，最大面积 $16900\pi\ \mathrm{m^2}$。设 $OM=m$，半径 $r=\dfrac{680-3m}5$，由 $r-m\geqslant80$ 与 $r-(60-m)\geqslant80$ 得 $10\leqslant m\leqslant35$，$m=10$ 时 $r=130$ 最大
  \item[\textbf{19.}] (1) 见解析版；(2) $\left(-\infty,-\dfrac13\right]$。不等式化为 $m\leqslant\dfrac{\mathrm e^{-x}-1}{\mathrm e^x+\mathrm e^{-x}-1}$ 对 $x>0$ 恒成立，右端最小值 $-\dfrac13$（$x=\ln2$ 时取得）；(3) 由条件得 $a>1$；当 $1<a<\mathrm e$ 时 $\mathrm e^{a-1}<a^{\mathrm e-1}$，当 $a=\mathrm e$ 时二者相等，当 $a>\mathrm e$ 时 $\mathrm e^{a-1}>a^{\mathrm e-1}$。依据：$f(x_0)>2$ 且 $-x_0^3+3x_0\leqslant2$，故 $a>1$；再由 $g(t)=\dfrac{\ln t}{t-1}$ 在 $(1,\mathrm e)$ 上大于 $\dfrac1{\mathrm e-1}$、在 $(\mathrm e,+\infty)$ 上小于 $\dfrac1{\mathrm e-1}$ 比较两幂
  \item[\textbf{20.}] (1) 见解析版；(2) $d=-1$。由 $S_2=2+d$ 必须是某项得 $m=2+\dfrac1d$ 为正整数，结合 $d<0$ 只能 $d=-1$，再验证 $m=\dfrac{n^2-3n+4}2$ 恒为正整数；(3) 见解析版
\end{itemize}

\gkpart{附加题}
\begin{itemize}[leftmargin=2.6em]
  \item[\textbf{21.}] (1) 见解析版；(2) $\dfrac72$。由 $A\boldsymbol\alpha=B\boldsymbol\alpha$ 得 $-2+2y=2+y$ 与 $2+xy=4-y$，解得 $y=4$，$x=-\dfrac12$；(3) $8\sqrt2$。直线为 $x+y-3=0$，与 $y^2=4x$ 交于 $(1,2)$、$(9,-6)$，故 $AB=8\sqrt2$；(4) 见解析版
  \item[\textbf{22.}] (1) $\dfrac5{18}$。同色取法 $\mathrm C_4^2+\mathrm C_3^2+\mathrm C_2^2=10$，任取 $2$ 球共 $\mathrm C_9^2=36$ 种；(2) $P(X=2)=\dfrac{11}{14}$，$P(X=3)=\dfrac{13}{63}$，$P(X=4)=\dfrac1{126}$，$E(X)=\dfrac{20}9$。按 $\mathrm C_9^4=126$ 计数：$X=4$ 仅四红 $1$ 种，$X=3$ 共 $26$ 种，其余 $99$ 种
  \item[\textbf{23.}] (1) $-1$。记 $g_n(x)=nf_{n-1}(x)+xf_n(x)$，可证 $g_n(x)=\cos\left(x+\dfrac{(n-1)\pi}2\right)$，故所求为 $g_2\left(\dfrac\pi2\right)=-\sin\dfrac\pi2$；(2) 见解析版
\end{itemize}
